\subsubsection*{4.18.}
先处理 $X$ 光滑且 $(X_0)_{\mathrm{red}}$ 是 $X$ 中的正规交叉除子的情形。由于各定理是局部的，可假设 $f$ 是从 $X=D^e$ 到 $D$ 的映射
\[
f(\underline z)=\prod_{i=1}^d z_i^{n_i},
\qquad 0<d\leq e,\quad n_i>0,
\]
这里 $D$ 表示标准圆盘 $D=\{z\in\mathbb C\mid |z|<1\}$。只需证明 (4.15.1)（分别为 (4.17.1)）在 $0\in D^e$ 处的纤维是拟同构。

设 $U$ 是 $D^e$ 中 $0$ 的形如 $|z_i|<a$ 的邻域，其中 $0<a\leq1$。在 (4.1.1) 中取 $U$ 的逆像，定义 $U_0$、$\bar U$、$U^*$ 和 $\bar U^*$。$U^*$ 或 $\bar U^*$ 的上同调不依赖于 $a$，因此
\[
H^*(\bar U^*,V)\xrightarrow{\ \sim\ }
  \bigl(R^*\Psi_\eta(V)\bigr)_0.
\]
$X^*=(D^*)^d\times D^{e-d}$ 的基本群 $\mathbb Z^d$（也是 $U^*$ 的基本群）为交换群。因此 $V$ 具有一个有限滤过，其相继商为秩一局部系统。通过逐步化简，可假设 $V$ 秩为一。

对 4.14 中的 $\widetilde{\mathcal V}$，用 $j_*^m\Omega_{X^*}^*(V)$ 表示 $j_*\Omega_{X^*}^*(V)$ 的子复形，其 $p$ 次项为 $(\Omega_X^p\otimes\widetilde{\mathcal V})^m$。由 [1], II 6.4.1 和 6.8，拟同构 $V\xrightarrow{\ \approx\ }\Omega_{X^*}^*(V)$（全纯 Poincaré 引理）诱导的自然映射
\begin{equation}\tag{4.18.1}
j_*^m\Omega_{X^*}^*(V)\longrightarrow Rj_*V,
\end{equation}
是拟同构。因此
\begin{equation}\tag{4.18.2}
H^*\bigl((j_*^m\Omega_{X^*}^*(V))_0\bigr)
 \xrightarrow{\ \sim\ }(R^*j_*V)_0
 \xleftarrow{\ \sim\ }H^*(U^*,V).
\end{equation}
\footnote{原印本 (4.18.2) 右端为 $H^*(U,V)$，但 $V$ 定义在 $U^*$ 上；英文恢复了显然应有的 $U^*$。}
这是证明的 4.18 部分唯一使用的解析结果。

\subsubsection*{引理 4.18.3.}
\underline{对秩一的 $V$，有}

\noindent\textup{(i)}\qquad
\uline{若 $V$ 非平凡：}\quad $H^*(U^*,V)=0$，

\medskip
\noindent\textup{(ii)}\qquad
\uline{若 $V=\underline{\mathbb C}$：}\quad
$H^*(U^*,V)=\bigwedge\nolimits^*\mathbb C^d$。

事实上，$U^*$ 具有环面的同伦型。

设 $\eta$ 是 $H^1(D^*,\mathbb Z)$ 的生成元，也表示其在 $H^1(U^*,\mathbb Z)$ 中的逆像。

\subsubsection*{引理 4.18.4.}
\underline{对秩一的 $V$，自然映射}
\[
H^*(U^*,V)
 \xrightarrow[(4.18.1)]{\ \sim\ }
H^*\bigl((j_*^m\Omega_{X^*}^*(V))_0\bigr)
 \longrightarrow
H^*\bigl(\Psi_\eta^{mu}(\Omega_{X^*}^*(V))_0\bigr)
\]
\underline{将 $H^*(\Psi_\eta^{mu}(\Omega_{X^*}^*(V))_0)$ 等同于 $H^*(U^*,V)/\eta H^*(U^*,V)$。}
\footnote{原印本左端为 $H^*(U^*,\mathbb C)$，但假设与结论涉及的都是一般的秩一系统 $V$；英文恢复为 $V$。}

使用 4.12(ii) 以及恒等式 $d\log f=\sum_i n_i\,dz_i/z_i$。复形 $\Psi_\eta^{mu}(\Omega_{X^*}^*(V))_0$ 等同于由双复形 $K$（使用直和）关联的单复形 $\mathrm{s}K$，其各项为
\[
K^{-p,n+p}=\frac{(\log f)^p}{p!}\cdot
  (j_*^m\Omega_{X^*}^n(V))_0.
\]
此复形位于一个不良象限 $(p\leq0,q\geq0,p+q\geq0)$ 中。不过，$\mathrm{s}K$ 是子复形 $\mathrm{s}\,\sigma_{p\geq p_0}(K)$ 的归纳极限，其各项为
\[
(\mathrm{s}\,\sigma_{p\geq p_0}K)^n
=\bigoplus_{\substack{p+q=n\\p\geq p_0}}K^{pq},
\]
这已足以使下面的论证成立。

按第一次数滤过的谱序列为
\[
E_1^{-p,n+p}=H^n(U^*,V)\qquad(p\geq0),
\]
其中
\[
d_1=2\pi i\,\eta\wedge:
H^n(U^*,V)\longrightarrow H^{n+1}(U^*,V).
\]
若 $V$ 非平凡，则由 4.18.3 得 $E_1=0$，而 4.18.4 归结为 $0=0$。假设 $V=\underline{\mathbb C}$。此时复形 $(E_1,d_1)$ 为可能截断的 Koszul 复形，由 Koszul 复形的无上同调性，得到
\[
E_2^{pq}=0\qquad\text{当 }p\ne0,
\]
及
\[
E_2^{0n}=H^n(U^*,V)/\eta H^{n-1}(U^*,V),
\]
这就证明了 4.18.4。

对 $\beta\in\mathbb C^*$，用 $V_\beta$ 表示 $D^*$ 上的秩一局部系统，其在 $\widetilde D^*$ 上的逆像已平凡化，而其单值为 $\beta$。其在 $U^*$ 上的逆像也使用相同记号。在 $\bar U^*$ 上，$V\otimes V_\beta$ 的逆像典范同构于 $V$ 的逆像。

\subsubsection*{引理 4.18.5.}
\underline{对秩一的 $V$，映射}
\[
\bigoplus_{\beta\in\mathbb C^*}
H^*(U^*,V\otimes V_\beta)/
 \eta H^*(U^*,V\otimes V_\beta)
\longrightarrow H^*(\bar U^*,V)
\]
\underline{是双射。}

设 $k$ 为各 $n_i$ 的最大公因数。空间 $\bar U^*$ 有 $k$ 个连通分支，以 $k$ 次单位根为指标。在以 $\xi$ 为指标的分支 $\bar U_\xi^*$ 上，$\widetilde D^*$ 上函数 $z^{1/k}$ 的逆像等于 $U^*$ 上函数
\[
\xi\prod_{i=1}^d z_i^{n_i/k}
\]
的逆像。分支 $\bar U_\xi^*$ 具有环面的同伦型，且序列
\[
0\longrightarrow\pi_1(\bar U_\xi^*)
 \longrightarrow\pi_1(U^*)\longrightarrow\pi_1(D^*)
\]
正合。用 $\varphi:\mathbb Z^d\to\mathbb Z$ 表示映射 $\pi_1(U^*)\to\pi_1(D^*)$。若 $\ker(\varphi)$ 在 $V$ 上作用非平凡，则 4.18.5 两端都为零。\footnote{原印本此处引用 4.18.4；所讨论的两端属于 4.18.5。}
否则，将 $V$ 换为适当的 $V\otimes V_\beta$ 后，可假设 $V$ 平凡，即 $V=\mathbb C$。此时在复上同调中有
\[
H^0(\bar U^*)\otimes
 H^*(U^*)/\eta H^*(U^*)
 \xrightarrow{\ \sim\ }H^*(\bar U^*).
\]
只需注意到：当 $\beta^k=1$ 时 $V_\beta$ 在 $U^*$ 上平凡，以及
\[
\bigoplus_{\beta^k=1}H^0(U^*,V_\beta)
=\bigoplus_{\beta^k=1}\mathbb C
 \xrightarrow{\ \sim\ }H^0(\bar U^*,\mathbb C),
\]
且当 $\beta^k\ne1$ 时 $H^*(U^*,V_\beta)=0$，便得结论。

\subsubsection*{注 4.18.6.}
若 $V$ 拟幺幂，4.18.5 中对 $\mathbb C^*$ 的求和可换为对单位根的求和；证明相同。

应用 4.12 得到交换图
\begin{center}
\begin{tikzcd}[
  column sep=2.2em,
  row sep=large,
  cells={nodes={font=\scriptsize}}
]
H^*\Psi_\eta^m(\Omega_{X^*}^*(V))_0
  \arrow[d]
& \displaystyle\bigoplus_\beta
  H^*\Psi_\eta^{mu}(\Omega_{X^*}^*(V\otimes V_\beta))
  \arrow[l,"\sim"',"4.12"]
& \displaystyle\bigoplus_\beta
  H^*(U^*,V\otimes V_\beta)/
    \eta H^*(U^*,V\otimes V_\beta)
  \arrow[l,"\sim"',"4.18.4"]
  \arrow[d,"\wr"',"4.18.5"]
\\
H^*R\Psi(V)_0
& {}
& H^*(\bar U^*,V)\arrow[ll,"\sim"']
\end{tikzcd}
\end{center}
由此得到 4.15。将 $\Psi^m$ 换为 $\Psi^{mqu}$，将对 $\beta$ 的求和换为对单位根的求和，并将 4.18.5 换为 4.18.6，即得命题 4.17。
